\documentclass[11pt,a4paper]{article}
\usepackage[utf8]{inputenc}
\usepackage[margin=1in]{geometry}
\usepackage{amsmath,amssymb}
\usepackage{booktabs}
\usepackage{hyperref}
\usepackage{enumitem}
\usepackage{microtype}
\usepackage{titlesec}

\titleformat{\section}{\large\bfseries}{\thesection}{1em}{}[\titlerule]
\setlist[itemize]{noitemsep, topsep=3pt, leftmargin=1.5em}
\setlist[enumerate]{noitemsep, topsep=3pt, leftmargin=1.5em}

\title{\textbf{Project Progress Report: Week 2\\
Low-Latency, Privacy-Preserving Dynamic ASL Recognition via Geometric Trajectory Modeling}}
\author{\textbf{Coordinate Trajectory Tracking \& Recognition Pipeline}}
\date{September 2026}

\begin{document}
\maketitle

\section*{Project Overview}
This project investigates low-latency, privacy-preserving American Sign Language (ASL) recognition using coordinate trajectory tracking from monocular webcam feeds. Rather than relying on computationally heavy 3D-CNNs operating directly on high-dimensional raw video frames, the proposed architecture decouples geometric feature extraction from temporal sequence modeling. We extract 21 3D hand keypoints per frame via Google MediaPipe Hands, normalize coordinates relative to the wrist and metacarpophalangeal (MCP) joints to achieve scale and position invariance, and pass the resulting sequences into a Bidirectional Gated Recurrent Unit (BiGRU) and a 1D Temporal Convolutional Network (TCN) to classify dynamic word-level signs.

\section*{Current Progress (Week 2)}
During Week 2, engineering focused on completing the spatial and temporal preprocessing stages, structuring the PyTorch tensor pipeline, and establishing benchmark baseline classifiers:

\begin{enumerate}
    \item \textbf{Spatial Normalization and Geometric Invariance Verification:} Implemented and verified coordinate normalization across all 21 hand landmarks. Coordinates are translationally centered relative to the wrist (Landmark 0) via:
    \[
        P'_i = P_i - P_{\text{wrist}} \quad \forall i \in \{0, \dots, 20\}
    \]
    Scale invariance is achieved by normalizing with respect to the Euclidean distance between the wrist and the middle-finger metacarpophalangeal (MCP) joint (Landmark 9):
    \[
        d_{\text{scale}} = \| P_{\text{mcp}} - P_{\text{wrist}} \|_2, \quad P''_i = \frac{P'_i}{d_{\text{scale}} + \epsilon}
    \]
    Empirical validation confirmed strict scale and position invariance, yielding a maximum residual discrepancy of $\Delta_{\max} < 8.16 \times 10^{-8}$ under arbitrary affine spatial perturbations.

    \item \textbf{PyTorch Dataset and DataLoader Construction:} Implemented a modular \texttt{ASLKeypointDataset} and multi-split \texttt{DataLoader} pipeline. Trajectories are resampled to a uniform temporal window of $T = 30$ frames (1.0 second at 30 FPS) via 1D spline interpolation, yielding tensor batches of dimension:
    \[
        (\text{batch\_size}, 30, 63)
    \]
    The dataset incorporates kinematic data augmentations for training robustness: Gaussian coordinate jittering ($\sigma = 0.01$), random isotropic scaling ($s \in [0.92, 1.08]$), and in-plane rotational perturbations ($\theta \in [-12^\circ, +12^\circ]$).

    \item \textbf{Kinematic Trajectory Feature Engineering:} To establish a classical baseline floor before training deep recurrent networks, we engineered 714 aggregated summary features across each 30-frame sequence. Features capture both static posture and dynamic trajectory kinematics:
    \begin{itemize}
        \item Spatial distribution: Coordinate mean, standard deviation, minimum, maximum, and bounding range ($5 \times 63$ features).
        \item Net trajectory displacement vector: $P_{T-1} - P_0$ ($63$ features).
        \item First-order velocity dynamics ($\Delta P / \Delta t$): Mean velocity, maximum instantaneous velocity, and velocity standard deviation ($3 \times 63$ features).
        \item Second-order acceleration dynamics ($\Delta^2 P / \Delta t^2$): Mean and maximum acceleration ($2 \times 63$ features).
        \item Cumulative 3D Euclidean path lengths per individual landmark joint ($21$ features).
    \end{itemize}

    \item \textbf{Non-Deep-Learning Baseline Benchmarking:} Trained tuned classical classifiers on the engineered 714-dimensional kinematic feature space over the 30-word target vocabulary:
    \begin{itemize}
        \item \textbf{Random Forest Classifier} ($n_{\text{estimators}} = 150$, balanced class weighting).
        \item \textbf{Support Vector Classifier (SVC)} with Radial Basis Function (RBF) kernel ($C = 3.0$, calibrated probability outputs, standard scaling pipeline).
    \end{itemize}
\end{enumerate}

\subsection*{Baseline Milestone Results}
Evaluation across held-out test instances yielded the following performance floor:

\begin{table}[h]
\centering
\begin{tabular}{lcccc}
\toprule
\textbf{Model Architecture} & \textbf{Top-1 Accuracy} & \textbf{Top-5 Accuracy} & \textbf{Macro F1-Score} & \textbf{Weighted F1-Score} \\
\midrule
Random Forest Classifier & \textbf{53.33\%} & \textbf{86.67\%} & \textbf{45.00\%} & \textbf{45.00\%} \\
Support Vector Classifier (SVC) & 40.00\% & 70.00\% & 31.89\% & 31.89\% \\
\bottomrule
\end{tabular}
\caption{Baseline classification metrics on the 30-word dynamic ASL vocabulary.}
\end{table}

The Random Forest model attained an 86.67\% Top-5 accuracy, confirming that coordinate trajectories combined with first- and second-order kinematic derivatives carry significant discriminative signal even prior to recurrent deep learning.

\section*{Challenges and Issues}
\begin{itemize}
    \item \textbf{Fine-Grained Geometric Minimal Pairs:} Certain signs (e.g., \textit{hello} vs. \textit{goodbye}, or \textit{food} vs. \textit{eat}) exhibit overlapping trajectory paths and differ primarily in subtle finger curling dynamics and inflection timing. While the Random Forest struggles with Top-1 disambiguation on these subtle pairs (limiting Top-1 to 53.33\%), the correct gloss consistently appears in the Top-5 predictions (86.67\%). This reinforces the need for deep temporal sequence models (BiGRU/TCN) to learn non-linear temporal phase dependencies.
    \item \textbf{Tracking Noise Sensitivity in Velocity Derivatives:} Discrete numerical differentiation of raw keypoints ($\Delta P / \Delta t$) amplifies high-frequency tracking jitter from the monocular detector. Applying temporal spline smoothing prior to velocity computation proved critical to avoid noisy acceleration spikes.
\end{itemize}

\section*{Next Steps (Week 3 Schedule)}
For Week 3, the schedule focuses on designing, training, and benchmarking the primary deep temporal sequence architectures:
\begin{itemize}
    \item \textbf{Bidirectional GRU Architecture:} Implement a 2-layer BiGRU with hidden dimension $H = 128$, dropout regularization ($p = 0.3$), and a temporal attention pooling layer operating over normalized input sequences $(\text{batch\_size}, 30, 63)$.
    \item \textbf{1D Temporal Convolutional Network (TCN):} Construct a dilated 1D-TCN backbone with exponentially increasing dilation factors ($d \in \{1, 2, 4, 8\}$) and residual skip connections to capture multi-scale temporal context with minimal inference latency.
    \item \textbf{Ablation Study:} Conduct a comparative evaluation between the Random Forest baseline, the BiGRU, and the 1D-TCN across Top-1 accuracy, Top-5 accuracy, inference latency (ms), and parameter count.
\end{itemize}

\section*{Questions or Feedback Addressed}
\begin{itemize}
    \item \textbf{Two-Hand vs. Single-Hand Representation:} We recommend adopting a \textbf{zero-padded 126-dimensional tensor representation ($42 \times 3$ features)} across the pipeline. In dynamic ASL, numerous essential signs within our target vocabulary (e.g., \textit{book}, \textit{family}, \textit{house}, \textit{more}, \textit{stop}) fundamentally require two interacting hands. Restricting the architecture to 63 single-hand dimensions causes severe semantic ambiguity. Zero-padding absent hands preserves uniform tensor shapes for batching without introducing spurious movement artifacts.
    \item \textbf{Evaluation Metrics (Top-1 vs. Top-5 \& Macro F1):} We recommend \textbf{reporting Top-5 accuracy alongside Macro-Averaged F1-score and Top-1 accuracy}. Given fine-grained sign similarities, Top-5 accuracy is the standard benchmark metric in sign language literature (e.g., WLASL, MS-ASL) and confirms whether candidate pools capture the intended gesture. Concurrently, Macro F1 ensures evaluation invariance against natural sign frequency imbalances.
\end{itemize}

\end{document}
